\chapter*{KẾT LUẬN VÀ ĐỀ NGHỊ}\addcontentsline{toc}{chapter}{Kết luận và đề nghị}
\section*{Kết luận}
Luận văn xây dựng mô hình gồm bảy yếu tố tác động đến quyết định chọn trường/ngành của học sinh THPT tại ĐBSCL: phù hợp cá nhân, triển vọng nghề nghiệp, danh tiếng--chất lượng, chi phí--hỗ trợ, ảnh hưởng xã hội, truyền thông tuyển sinh và vị trí--điều kiện học tập. Mô hình tích hợp góc nhìn hành vi, đầu tư giáo dục và phù hợp người--môi trường, đồng thời cung cấp thang đo 32 biến quan sát và quy trình kiểm định đầy đủ.

Kết quả minh họa trên 486 quan sát cho thấy cả bảy giả thuyết có chiều dương; FIT và CAR nổi bật nhất. Tuy nhiên, đây chỉ là mẫu trình bày, không phải bằng chứng thực nghiệm. Đóng góp thực sự của bản nộp cuối phụ thuộc vào việc thu thập dữ liệu đúng đạo đức, phân tích tái lập được và diễn giải thận trọng.

\section*{Đề nghị}
ĐHCT và các trường đại học trong vùng nên ưu tiên tư vấn phù hợp ngành, minh bạch triển vọng việc làm và chi phí, đồng thời cá nhân hóa thông tin theo địa bàn. Trường THPT cần đưa trải nghiệm nghề vào hướng nghiệp; gia đình nên hỗ trợ thay vì quyết định thay. Cơ quan quản lý có thể phát triển cơ sở dữ liệu ngành học--việc làm liên thông, chuẩn hóa cách công bố tỷ lệ việc làm.

\section*{Hạn chế và hướng nghiên cứu tiếp theo}
Thiết kế cắt ngang không xác lập nhân quả; tự báo cáo có thể chịu thiên lệch xã hội; lấy mẫu theo trường có hiệu ứng cụm; ý định không đồng nhất với hành vi nhập học. Nghiên cứu tiếp theo nên theo dõi học sinh từ lúc đăng ký đến sau nhập học, dùng mô hình đa mức theo trường/tỉnh, kiểm định tính bất biến thang đo giữa nhóm, kết hợp phỏng vấn sâu và đối chiếu dữ liệu nguyện vọng thực tế. Có thể tách DEC thành quyết định chọn trường và chọn ngành để kiểm định mô hình SEM.

